\documentclass[11pt,a4paper]{ltjsarticle}
\usepackage[haranoaji]{luatexja-preset}
\usepackage{lmodern,amsmath,amssymb,bm,geometry,booktabs,tabularx,array,enumitem,xcolor,hyperref}
\geometry{margin=22mm}
\hypersetup{colorlinks=true,linkcolor=blue!40!black,urlcolor=blue!45!black,pdftitle={実験III v0.1 実施報告},pdfauthor={Toru Takahashi}}
\setlist[itemize]{leftmargin=1.7em,itemsep=2pt,topsep=3pt}
\setlist[enumerate]{leftmargin=2.0em,itemsep=3pt,topsep=3pt}
\setlength{\parskip}{3pt}\setlength{\parindent}{1em}\renewcommand{\arraystretch}{1.16}
\setlength{\emergencystretch}{2em}\raggedbottom\pagestyle{plain}
\definecolor{navy}{HTML}{174A6E}\definecolor{teal}{HTML}{187D7A}\definecolor{pale}{HTML}{EEF5F7}\definecolor{pass}{HTML}{176B55}
\newcommand{\PASS}{\textcolor{pass}{\textbf{PASS}}}
\newcommand{\FAIL}{\textcolor{red!65!black}{\textbf{FAIL}}}

\title{\vspace{-1.3em}\textbf{実験III v0.1 実施報告}\\\large 世界形成の差を生む計算要因の分離と多変数閉ループ}
\author{Toru Takahashi\quad 廣安研究室・独立非常勤研究員}
\date{2026年8月24日}

\begin{document}
\maketitle

\begin{center}
\fcolorbox{teal}{pale}{\parbox{0.91\textwidth}{
\textbf{結論。}有限状態の資源管理環境と小規模な線形共有表現を用い、課題枠組み、因果モデル、予測地平、評価地平、リスク感度、表現容量、学習機会、課題枠組み可塑性、代替情報精度、停止規則を個別に操作できる測定系を構成した。さらに、これらの状態を外生的な更新則で閉ループ化すると、履歴に応じた二状態、時間尺度分離、有限摂動による転換、別表層環境での状態維持が得られた。一方、候補座標$q_t=\kappa_t$のプール予測連関は摂動量を条件づけると保たれず、転換を個体差レベルで記述する低次元反応座標は同定できなかった。更新則そのものは学習されておらず、人間の志向性や類型を実証した結果ではない。
}}
\end{center}

\section{目的}

実験Iでは、取得しなかった情報と取得後に圧縮された情報を分離し、有限な観測・表現資源が世界表現を分岐させうることを確認した。実験IIでは、現在の表現が次の学習対象を偏らせる閉ループにより、履歴依存の二状態と摂動後の復帰・転換、ヒステリシスを構成した。

実験IIIの目的は、人間的な類型を人工系へ当てはめることではない。結論や方策の差が、\textbf{何を問題とするか、何が起こると予測するか、どこまでを評価するか、どの資源が不足するか、いつ推論を止めるか}のどこから生じたのかを、介入によって分離することである。そのうえで、複数の状態が閉ループ内で準安定化し、将来の状態遷移を予測する低次元量へ縮約できるかを調べた。

\section{共通環境と実装}

\subsection{有限状態の資源管理環境}

状態は、蓄積量 $x\in\{0,\ldots,4\}$、需要状態 $d\in\{0,1\}$、設備状態 $h\in\{0,1\}$からなる20状態である。行為は供給量と保守の組で、競合環境では5種類、帰無環境では現在供給と保守を同時最大化できる行為を加えた6種類とした。主要帰結は現在不足と将来障害リスクである。地平10の有限状態動的計画法により、生成過程を知るOracle方策を厳密に求めた。

通常観測は需要状態と設備状態を正答率0.62で示す。需要センサと設備センサは対応状態を正答率0.94で示し、非対応状態の情報を含まない。観測予算と、後続推論へ保持できる表現スロット数を別に設定した。

\subsection{共有表現の学習系}

資源形成実験では、独立な二信号を回転混合した入力から需要信号と設備信号を復号する線形共有表現を自作した。最大ボトルネックは2次元で、1次元条件では第2表現成分を使えない。各対象の正規化復号誤差を $D_d,D_h$ とし、表現偏差を
\begin{equation}
s=\frac{D_h-D_d}{D_h+D_d}
\end{equation}
とした。$s>0$は需要側、$s<0$は設備側への相対的な特化を表す。安定化型選択則は現在の$s$と遅い課題枠組み状態から次の学習対象を確率的に選ぶ。

ニューラルネットワーク用ライブラリには依存せず、NumPyによる有限状態Oracleと勾配更新を実装した。すべてCPUで完了する。

\section{III-A0：環境と測定量の校正}

遷移確率、費用、Oracle、帰結競合、センサ選択性、資源水準を結果を見る前の五仮説群で検査した。全202個のOracle重み条件、12個の資源条件、4個のセンサ条件が揃い、完全性検査を通過した。

\begin{center}\small
\begin{tabular}{lrrr}
\toprule
現在不足の重み & Oracle目的値 & 一様ランダム & 相対改善\\\midrule
0.25 & 0.090477 & 0.139441 & 35.1\%\\
0.50 & 0.176332 & 0.240146 & 26.6\%\\
0.75 & 0.261074 & 0.340852 & 23.4\%\\\bottomrule
\end{tabular}
\end{center}

競合環境には73個の非劣点があり、現在不足の範囲は0.398179、障害リスクの範囲は0.155361だった。帰無環境には両帰結を同時に最小化する共通方策が存在した。需要・設備センサは対応対象の正答率を0.62から0.94へ改善し、非対象の正答率を変えなかった。観測予算2・表現スロット2で平均Bayes誤りは0.06となり、それ以上の資源による改善はなかった。A0の五仮説群はすべて\PASS となった。

\section{III-A1：課題、予測、評価の一要因分離}

\subsection{最初の確認と設計監査}

課題枠組み、因果モデル、予測地平、評価地平、リスク感度を一つずつ変えた最初の確認では、課題枠組みと予測地平だけが成立し、因果モデル、評価地平、リスク感度は\FAIL となった。因果モデル介入は使用による設備劣化だけを外しており、方策不一致率0.025、予測誤差0.000425と操作が弱かった。評価地平によるリスク差は0.003901で事前閾値0.005に届かなかった。リスク感度の絶対低下閾値0.02は、基準危険状態率0.018488より大きく、達成不能だった。

この不成立を残したうえで、介入強度を探索的に校正した。直接効果しか扱わない因果モデル、評価割引0.10、相対リスク低下を含む判定を固定し、未使用の40環境で確認した。

\begin{center}\small
\begin{tabularx}{0.94\textwidth}{Xrrr}
\toprule
測定量 & 平均 & 95\%下限 & 95\%上限\\\midrule
課題対象の正答率 & 0.9416 & 0.9364 & 0.9469\\
因果モデルの方策不一致率 & 0.1369 & 0.1300 & 0.1438\\
因果モデル誤指定の予測誤差 & 0.1570 & 0.1535 & 0.1607\\
予測地平による目的改善 & 0.0080 & 0.0077 & 0.0082\\
評価地平による障害リスク改善 & 0.0123 & 0.0119 & 0.0128\\
リスク感度による危険状態の相対低下 & 0.6955 & 0.6825 & 0.7088\\\bottomrule
\end{tabularx}
\end{center}

独立確認では五仮説群がすべて\PASS となった。遠い帰結を予測できることと、それを評価へ含めることは、別々に方策を変える要因として分離できた。

\section{III-A2：有限資源の形成効果と維持効果}

形成期と測定期について、非拘束Uと拘束Cを2$\times$2にした。非拘束Uは2次元表現かつ両対象更新、拘束Cは1次元表現かつ一対象更新である。

\begin{center}\small
\begin{tabular}{lrrr}
\toprule
測定量 & 平均$|s|$ & 95\%下限 & 95\%上限\\\midrule
形成期C終了 & 0.970607 & 0.969614 & 0.971610\\
形成期U終了 & 0.054956 & 0.040151 & 0.071835\\
UU終端 & 0.059168 & 0.043355 & 0.077692\\
UC終端 & 0.986192 & 0.985659 & 0.986713\\
CU終端 & 0.130342 & 0.105671 & 0.157305\\
CC終端 & 0.985647 & 0.985167 & 0.986134\\\bottomrule
\end{tabular}
\end{center}

拘束下では形成時・測定時とも強い分岐が生じた。形成後に資源を解放したCUでは、$|s|$が平均0.8403低下し、拘束下で形成された枝は同じ強度では維持されなかった。ただしCUの$|s|=0.1303$はUUの0.0592より大きく、形成履歴の偏差が完全に消失したわけではない。資源解放は制御条件の変更であるため、この結果は、履歴がどの枝へ入るかを決めることと、その枝を支える力学条件が別であることを示す。この構成系では、継続する資源拘束が強い枝構造の形成と維持の双方に関与した。

続いて表現容量と学習機会を分離した。容量1で両対象を同時更新しても$|s|=0.9037$となり、対称な圧縮解の縮退から一方への特化が生じた。容量2でも一対象更新では$|s|=0.3368$、容量2・両対象更新では0.0659だった。表現容量不足による復号誤差和の増分は0.9886、学習機会選択による偏差増分は0.2709だった。両者は別々の分岐源である。

\section{III-B：課題再編と停止規則}

\subsection{代替情報精度・可塑性・二つの地平}

潜在設備状態について、代替情報精度、課題枠組み可塑性、予測地平、評価地平を2$\times$2$\times$2$\times$2にした。現在の課題枠組みが利用できる状態を制限し、代替情報精度は代替センサの正答率、可塑性は枠組み変更の可否として実装した。課題再編則自体は外部設計である。

\begin{center}\small
\begin{tabular}{lrrr}
\toprule
対応差 & 平均 & 95\%下限 & 95\%上限\\\midrule
長予測・長評価での切替率 & 0.2978 & 0.2924 & 0.3034\\
二つの地平の相互作用 & 0.4125 & 0.4058 & 0.4190\\
代替情報精度による状態一致率の改善 & 0.4085 & 0.3975 & 0.4196\\
代替情報精度による障害リスク改善 & 0.0187 & 0.0178 & 0.0197\\
可塑性による障害リスク改善 & 0.1119 & 0.1084 & 0.1154\\\bottomrule
\end{tabular}
\end{center}

長い予測地平だけ、または長い評価地平だけでは課題は再編されず、両方が長く、かつ枠組みが可塑的な場合に切替が生じた。これは、本環境と外部設計した再編則において、予測可能性、評価範囲、枠組み可塑性が連言的に作用した結果であり、一般的な必要条件を示すものではない。代替情報の精度を高めると切替回数を単純に増やすのではなく、課題枠組みと真の設備状態の一致率が改善した。この操作は代替情報源の品質を変えたものであり、代替観測を選ぶ傾向としての探索方策は操作していない。

\subsection{停止規則}

固定12観測、事後エントロピー閾値、限界情報利得閾値を比較した。停止後は新しい証拠を取得せず、全条件で総計算回数を30へ揃えた。

精度閾値規則は易しい試行より難しい試行へ平均18.93回多く観測を配分し、固定規則とほぼ同じ平均観測数で正答率を0.0463改善した。情報利得規則は正答率を0.0060改善しながら観測を6.15回削減した。停止規則の効果は総計算時間の差では説明できない。

\section{III-C：多変数閉ループと反応座標の監査}

課題枠組みを遅い状態$\kappa_t$、評価感度をさらに遅い状態として表現学習へ結合した。需要履歴と設備履歴の二腕を作った後、外部条件と更新則を同じにした。パイロットで候補座標を$q_t=\kappa_t$へ固定し、確認データでは変更しなかった。

80腕すべてが履歴に対応する枝を形成し、共通条件終了時の平均$|s|$は0.9854だった。8000更新の無摂動監視で自発転換は0、弱摂動後の復帰時間90\%分位は200更新であり、観測された時間尺度比の下限は40だった。表面の$s$は摂動4000--5600で平均差0.0012しかなかったが、転換率は0.4125変化した。入力回転を変えた同型環境でも枝維持率は1.000、終端aligned $s$は0.9898だった。

最初の確認ではK=3200を復帰優勢とする上限0.10に対し転換率0.150となり、有限転換域の仮説群だけが\FAIL となった。他の四仮説群は成立した。K=3200の不成立を維持し、未使用の別40 seed・80腕で弱水準K=1600と強水準K=6400を再確認した。

\begin{center}\small
\begin{tabular}{crrrrrrr}
\toprule
K & 0 & 1600 & 3200 & 4000 & 4800 & 5600 & 6400\\\midrule
転換率 & .000 & .000 & .138 & .450 & .688 & .838 & .938\\\bottomrule
\end{tabular}
\end{center}

用量反応は単調で、弱水準と強水準の差は0.938となり、有限の復帰域と転換域が独立集合で確認された。

\subsection{候補座標\texorpdfstring{$q_t=\kappa_t$}{q = kappa}の事後監査}

事前に固定したプール解析では、$q_t$の転換AUCは0.712、表面指標$s$は0.600だった。しかし、摂動量$K$の三水準をまとめた解析では、$q_t$が摂動強度に伴う変化を利用している可能性を除けない。この点を結果取得後の監査として、$K$水準内と$K$を統制した予測モデルで再評価した。

\begin{center}\small
\begin{tabular}{rrrrr}
\toprule
$K$ & 腕数 & 転換率 & AUC($\kappa$) & AUC($s$)\\\midrule
4000 & 80 & 0.463 & 0.565 & 0.537\\
4800 & 80 & 0.625 & 0.567 & 0.464\\
5600 & 80 & 0.875 & 0.380 & 0.474\\
プール & 240 & 0.654 & 0.712 & 0.600\\\bottomrule
\end{tabular}
\end{center}

$K$水準内では$\kappa_t$の予測連関は保たれなかった。さらにseedを分離した5分割交差検証で、$K$をカテゴリ変数として統制したロジスティックモデルを比較した。

\begin{center}\small
\begin{tabular}{lrrr}
\toprule
予測変数 & AUC & log loss & Brier score\\\midrule
$K+s$ & 0.6527 & 0.6001 & 0.2074\\
$K+s+\kappa$ & 0.6576 & 0.6004 & 0.2072\\\bottomrule
\end{tabular}
\end{center}

$\kappa_t$追加後のAUC差は0.0049にとどまり、log lossはわずかに悪化した。したがって、プール解析で得られた予測連関は$K$を条件づけると保たれず、$q_t=\kappa_t$を個体差レベルの転換反応座標として支持できない。$\kappa_t$が主として$K$の代理であった可能性は高いが、この監査だけで代理関係そのものを証明したわけではない。準安定な二状態は構成できた一方、その将来遷移を十分に記述する低次元反応座標は未同定である。

\section{統合して言えること}

\begin{enumerate}
\item 課題枠組み、因果モデル、予測地平、評価地平、リスク感度は、同じ有限状態環境で異なる診断パターンをもつ介入として分離できる。
\item 表現容量不足と、一度に学べる対象数の制限は、別々に表現偏差を生みうる。
\item この構成系では、継続する資源拘束が強い偏りの形成と維持の双方に関与した。拘束解放後には偏りが大きく縮小したが、弱い履歴差は残った。
\item 本環境と再編則では、長い予測地平、長い評価地平、課題枠組み可塑性が揃ったときに課題再編が生じた。三者は同一の性質ではない。
\item 代替情報精度、課題枠組み可塑性、二つの地平、停止規則は異なる診断量へ作用した。ただし探索方策の独立効果は未検証である。
\item 外部設計した多変数閉ループにおいて、時間尺度分離、有限の復帰域・転換域、別表層環境での枝維持を示す二状態を構成できた。一方、その転換を個体差レベルで予測する低次元反応座標は同定できなかった。
\end{enumerate}

\enlargethispage{2\baselineskip}
実験III v0.1は、世界形成差の計算要因を因果的に分離し、外生的更新則のもとで再結合する最小基盤を得た。準安定レジームの構成と反応座標の同定は別の課題であり、後者は未達である。

\section{限界と次段階}

\begin{itemize}
\item 課題再編、対象選択、遅い状態の更新則は外部から設計されている。規則そのものの獲得は実験Vの対象である。
\item $q_t=\kappa_t$のプール予測連関は$K$を条件づけると保たれず、転換反応座標として支持されなかった。表面指標$s$だけでも十分でなく、必要な座標の次元と内容は未同定である。
\item 代替センサ精度の操作は情報品質を変えたものであり、探索方策そのものの介入ではない。
\item 別環境は入力回転を変えた同型課題であり、意味内容の異なる課題横断性ではない。
\item 線形共有表現、二対象、単一の資源管理環境での結果である。人間の志向性、価値観、文化差への一般化はできない。
\item 自発転換0は滞在時間が無限であることを意味しない。観測時間内の遷移率上限として扱う。
\end{itemize}

次段階では、代替情報源の精度を固定したまま、代替観測の選択確率または情報利得への重みだけを操作し、探索方策と課題枠組み可塑性を分離する。反応座標については$\kappa_t$の追加確認を優先せず、軌道データの遅い自由度の数、局所勾配応答、表現幾何、重み履歴などを含め、低次元座標が存在するかという問いから再検討する。その後、意味内容の異なる因果環境と、心理状態・構造関係、現前結合・地平、具象差異・抽象関係を個別に操作する課題へ進む。

\section{再現性}

パイロットと確認実験のseed集合を分離し、確認実験ごとに事前登録文、コードハッシュ、完全性検査、CSV、JSONを保存した。初回確認で不成立となった仮説群と、設計上の手戻りを結果から除外していない。$q_t$の条件付き監査は結果取得後の事後解析であり、事前登録した判定とは区別して報告した。実行はNumPyのみで、GPUを必要としない。

主要な確認seedは、A1が3200--3239、A2が3400--3439、A2bが3500--3539、B1が3700--3739、B2が3900--3939、Cが4100--4139、Cの転換域再確認が4200--4239である。

\end{document}
